\documentclass[journal, a4paper]{IEEEtran}
\usepackage[utf8]{inputenc}
\usepackage[spanish]{babel}
\usepackage{amsmath}
\usepackage{graphicx}
\usepackage{pgfplots}
\usepackage{booktabs}
\usepackage{siunitx}

\pgfplotsset{compat=1.18}

\begin{document}

\title{Pre-informe: Caracterización y Análisis Transitorio del Diodo 1N4004}

\author{
    Juan José Tobar Álvarez, \textit{jutobara@unal.edu.co}\\
    Estudiante 2, \textit{correo2@e-mail.com}\\
    Estudiante 3, \textit{correo3@e-mail.com}
}

\maketitle

\section{Introducción y Justificación}
Este pre-informe caracteriza el diodo rectificador 1N4004, determinando su relación tensión-corriente, tiempo de recuperación inversa y corriente de saturación inversa \cite{guia}, a partir de la ecuación de Shockley $I_{d}=I_{s}(e^{V_d/\eta V_{T}})$, donde $I_s$ es la corriente de saturación inversa, $\eta$ el coeficiente de emisión y $V_T$ el voltaje térmico \cite{guia}.

Esta caracterización es relevante en entornos de alta exigencia: en electrónica espacial, los rectificadores calificados para satélites deben operar de forma confiable hasta 175 °C de temperatura de unión, lo que exige conocer con precisión el desplazamiento de $V_t$ e $I_s$ con la temperatura \cite{strad}; en el sector médico, el diseño de desfibriladores externos requiere diodos de recuperación rápida para controlar la descarga de energía y evitar corrientes inversas indeseadas \cite{allpcb}.

\section{Cálculos}
Para $I_{max}=50$ mA con $V_{fuente}=25$ V y asumiendo $V_d\approx0.7$ V:
$$R_{1} = \frac{25 - 0.7}{0.05} = 486~\Omega \;\approx\; 500~\Omega \text{ (valor de montaje)}$$

\section{Simulación}

\subsection{Efecto de la Temperatura (4.3.1)}
Se simuló el diodo a -25°C, 25°C y 75°C.

\begin{figure}[htbp]
    \centering
    \pgfplotstableset{comment chars={S}}
 
    \begin{tikzpicture}
        \begin{axis}[
            title={Curva Característica V-I a diferentes Temperaturas},
            xlabel={Voltaje en el diodo $V_{D}$ (mV)},
            ylabel={Corriente I(D1) (mA)},
            grid=major,
            width=\columnwidth,
            height=6cm,
            xmin=0, xmax=780,
            ymin=0, ymax=54,
            xtick distance=70,
            ytick distance=8,
            legend pos=north west,
            scaled x ticks=false,
            scaled y ticks=false
        ]
            % Curva 1: Temp = -25 (Color Azul)
            \addplot[color=blue, thick, restrict expr to domain={\coordindex}{0:250}]
                table[x expr=\thisrow{V(n002)}*1000, y expr=\thisrow{I(D1)}*1000] {Grafica1.txt};
            \addlegendentry{-25$^\circ$C}
 
            % Curva 2: Temp = 25 (Color Verde)
            \addplot[color=green, thick, restrict expr to domain={\coordindex}{251:501}]
                table[x expr=\thisrow{V(n002)}*1000, y expr=\thisrow{I(D1)}*1000] {Grafica1.txt};
            \addlegendentry{25$^\circ$C}
 
            % Curva 3: Temp = 75 (Color Rojo)
            \addplot[color=red, thick, restrict expr to domain={\coordindex}{502:752}]
                table[x expr=\thisrow{V(n002)}*1000, y expr=\thisrow{I(D1)}*1000] {Grafica1.txt};
            \addlegendentry{75$^\circ$C}
 
        \end{axis}
    \end{tikzpicture}
    \caption{Efecto de la temperatura en la curva del diodo 1N4004. Datos de simulación exportados desde LTSpice \cite{grafica1}.}
    \label{fig:temp_variation}
\end{figure}

Como muestra la Fig. \ref{fig:temp_variation}, al aumentar $T$ crecen $V_T$ e $I_s$ \cite{guia}, reduciendo la tensión umbral necesaria para conducir. Frente al modelo ideal ($I_d>0,V_d=0$) \cite{guia}, a baja temperatura el diodo se \emph{aleja} de la idealidad en la caída de tensión (requiere más $V_d$), pero se \emph{acerca} en la nitidez de la transición corte-conducción (curva más abrupta) y en la menor fuga inversa. $I_s$ es el parámetro con dependencia térmica dominante; $\eta$, asociado al mecanismo de conducción, debería variar poco con $T$.

\begin{table}[h]
    \centering
    \caption{$V_d$ a $I_d\approx50$ mA según temperatura}
    \begin{tabular}{ccc}
        \toprule
        \textbf{T} & \textbf{$V_d$ (V)} & \textbf{$I_d$ (mA)} \\
        \midrule
        -25°C & 0.7597 & 49.88 \\
        25°C & 0.6751 & 50.05 \\
        75°C & 0.5884 & 50.23 \\
        \bottomrule
    \end{tabular}
    \label{tab:temp_vd}
\end{table}
El coeficiente térmico resultante, $\Delta V_d/\Delta T\approx-1.71$ mV/°C, es cercano al valor típico de $-2$ mV/°C para uniones de silicio.

\subsection{Modelamiento del Diodo (4.3.2)}
$R_d$ y $V_t$ \cite{guia} se estiman con la secante entre el punto de máxima corriente y un punto de baja corriente de la curva a 25°C (Tabla \ref{tab:parametros}).

\begin{table}[h]
    \centering
    \caption{Puntos de la curva a 25°C \cite{grafica1}}
    \begin{tabular}{lcc}
        \toprule
        \textbf{Punto} & \textbf{$V_d$ (V)} & \textbf{$I_d$ (mA)} \\
        \midrule
        Bajo ($V_1=5$ V) & 0.5868 & 9.08 \\
        Medio ($V_1=12.5$ V) & 0.6380 & 24.41 \\
        Máximo ($V_1=25$ V) & 0.6751 & 50.05 \\
        \bottomrule
    \end{tabular}
    \label{tab:parametros}
\end{table}

Con Bajo-Máximo: $R_d\approx2.16~\Omega$, $V_t\approx0.567$ V. Con Bajo-Medio (en vez del máximo): $R_d'\approx3.34~\Omega$, $V_t'\approx0.557$ V. Es decir, $R_d$ \emph{aumenta} y $V_t$ \emph{disminuye} levemente al usar el punto medio, por la curvatura exponencial de la característica.

\subsection{Tiempo de Recuperación Inversa (4.3.3)}
Con onda cuadrada de 10 V, $t_{rr}=t_s+t_f$ \cite{guia}.

\begin{figure}[htbp]
    \centering
    \begin{tikzpicture}
        \begin{axis}[
            title={Respuesta Transitoria: Corriente vs Tiempo},
            xlabel={Tiempo ($\mu$s)},
            ylabel={I(D1) (mA)},
            grid=both,
            width=\columnwidth,
            height=5.5cm,
            legend pos=south east,
            scaled x ticks=false,
            scaled y ticks=false,
            xtick distance=3,
            ytick distance=3,
            tick label style={
                /pgf/number format/fixed,
                /pgf/number format/precision=1,
                /pgf/number format/fixed zerofill
            }
        ]
            \addplot[color=red, thick] table[
                x expr=\thisrow{time}*1e6,
                y expr=\thisrow{I(D1)}*1000
            ] {Grafica2.txt};
            \addlegendentry{Corriente}
        \end{axis}
    \end{tikzpicture}
    \caption{Tiempo de recuperación inversa para el diodo 1N4004. Datos de simulación exportados desde LTSpice \cite{grafica2}.}
    \label{fig:trr}
\end{figure}

Teóricamente $t_{rr}$ no debería depender de la frecuencia, pues obedece a la recombinación de portadores de la unión P-N; discrepancias a 100 kHz frente a la hoja de datos se atribuyen a aproximaciones numéricas del simulador. Al cambiar $R_1$ a 500 $\Omega$, $I_R$ aumenta y $t_s$ disminuye. En SPICE, $t_s$ y $t_f$ están gobernados por el tiempo de tránsito ($TT$) y la capacitancia de unión ($CJ0$), no por $I_s$ ni $\eta$ \cite{guia}.

\section{Montaje}
% Sección dejada en blanco de acuerdo con las instrucciones del pre-informe.

\section{Análisis de Resultados}
El diodo 1N4004 responde coherentemente con la ecuación de Shockley \cite{guia}: la temperatura desplaza la rodilla de conducción, y el transitorio evidencia $t_s$ y $t_f$ al forzar el corte \cite{guia}, confirmando que rectificadores de propósito general como el 1N4004 son más lentos que diodos de conmutación rápida (1N4148).
% Comparación con datos experimentales pendiente por falta de mediciones de laboratorio.

\section{Conclusiones}
\begin{itemize}
    \item La temperatura afecta inversamente a $V_t$ y directamente a $I_s$, con un coeficiente medido de $-1.71$ mV/°C, coherente con la teoría.
    \item $R_d$ y $V_t$ dependen del tramo de la curva usado para la secante: el punto medio da mayor $R_d$ y menor $V_t$ que el punto máximo, por la no linealidad exponencial del diodo.
    \item $t_{rr}$ depende de la recombinación del semiconductor, no de la frecuencia de excitación; sin embargo, $R_1$ sí modifica $t_s$ y $t_f$ al alterar la amplitud de la corriente de barrido.
\end{itemize}

\section{Referencias}
\begin{thebibliography}{1}
\bibitem{guia} Universidad Nacional de Colombia, \textit{Práctica 1: Caracterización del Diodo}, Electrónica Análoga I, 2026-II.
\bibitem{strad} STMicroelectronics, \textit{Rad-hard Diodes and Rectifiers}. [En línea]. Disponible en: https://www.st.com/en/space-products/rad-hard-diodes-and-rectifiers.html
\bibitem{allpcb} AllPCB, \textit{Defibrillator PCB Design for Beginners: A Step-by-Step Guide}, 2025. [En línea]. Disponible en: https://www.allpcb.com/blog/pcb-knowledge/defibrillator-pcb-design-for-beginners-a-step-by-step-guide.html
\bibitem{grafica1} LTSpice Simulation Data Export, \textit{Grafica1.txt}.
\bibitem{grafica2} LTSpice Simulation Data Export, \textit{Grafica2.txt}.
\end{thebibliography}

\end{document}